% Generated draft: requires mathematical and source review.
\documentclass[11pt]{article}
\usepackage[margin=25mm]{geometry}
\usepackage{fontspec}
\setmainfont{Latin Modern Roman}
\usepackage{amsmath,amssymb,mathtools,unicode-math}
\setmathfont{Latin Modern Math}
\usepackage{enumitem,needspace,xurl,hyperref}
\setlength{\parindent}{0pt}
\setlength{\parskip}{5pt}
\setlength{\emergencystretch}{3em}
\newcommand{\N}{\mathbb{N}}
\newcommand{\Z}{\mathbb{Z}}
\newcommand{\Q}{\mathbb{Q}}
\newcommand{\R}{\mathbb{R}}
\newcommand{\C}{\mathbb{C}}
\newcommand{\F}{\mathbb{F}}
\newcommand{\A}{\mathbb{A}}
\newcommand{\PP}{\mathbb{P}}
\newcommand{\E}{\mathbb{E}}
\newcommand{\eps}{\varepsilon}
\newcommand{\dd}{\,\mathrm{d}}
\begin{document}
\section*{UnsolvedMath ID: 3000083}
\textbf{Source problem code:} catalogue entry
\subsection{Definitions and mathematical statement}
For graphs \(G=(V,E)\) and \(H\), define
\[
\mathcal F_H(G)=\{F\subseteq E:(V,F)\text{ contains no subgraph isomorphic to }H\}.
\]
Subgraph here is not induced subgraph. An \(H\)-free edge set \(F\) is strongly maximal when
\[
\forall I\in\mathcal F_H(G),\qquad |I\setminus F|\le|F\setminus I|.
\]
For infinite sets these are cardinal comparisons of the two differences; they are not subtraction of infinite cardinalities and not merely a finite-exchange condition. The classification problem is to identify exactly the graphs \(H\) with
\[
\forall G\quad\exists F\in\mathcal F_H(G)\quad
\forall I\in\mathcal F_H(G):|I\setminus F|\le|F\setminus I|.
\]
The intended graph categories and cardinality ranges must be specified. For finite \(G\), whenever \(\mathcal F_H(G)\) is nonempty, an \(H\)-free set of largest cardinality already has this property, so the substantive general question involves infinite host graphs. Whether \(H\) itself is required finite, has an edge, or may have isolated vertices is not stated. Edgeless forbidden graphs may make the feasible family empty for some hosts and cannot silently be included or excluded.

When \(H\) is the path of \emph{length two}, with two edges and three vertices, its free edge sets are matchings in the usual simple-graph convention. The source states that every host graph has a strongly maximal matching. Its research summary uses the symbol \(P_2\) for this case, which differs from the convention that indexes a path by its number of vertices; the two-edge description is retained to avoid that ambiguity.

The supplied partial-progress summary says no full characterization was found, claims some other forbidden graphs fail the property without identifying them, and mentions strongly maximal matching results in graphs and finite-rank hypergraphs. These are supplied contextual claims, not a classification established here.
\subsection{Short English statement}
For which forbidden graphs does every host graph admit a free spanning subgraph that cannot gain more edges than it loses in any free replacement? The infinite-graph scope and forbidden-graph conventions are essential.
\subsection{Sources}
\begin{itemize}
\item Local UnsolvedMath registry, permanent ID \(3000083\), exact source code \texttt{catalogue entry}, ``Strongly maximal H-free spanning subgraph.''
\item Egres Open, \emph{List of open problems}, page of the same title; \url{https://oldlemon.cs.elte.hu/egres/open/List_of_open_problems}. HTML extraction; recorded access 2026-07-29.
\item The statement cites an unnamed source, page 16, Theorem 5.6, for strongly maximal matchings. The research summary names Aharoni--Berger--Ziv, Friedman, Erd\H{o}s--Pado, and Aharoni without complete citations or identified counterexamples for other \(H\).
\end{itemize}
\end{document}
